\IfFileExists{stacks-project.cls}{%
\documentclass{stacks-project}
}{%
\documentclass{amsart}
}

% For dealing with references we use the comment environment
\usepackage{verbatim}
\newenvironment{reference}{\comment}{\endcomment}
%\newenvironment{reference}{}{}
\newenvironment{slogan}{\comment}{\endcomment}
\newenvironment{history}{\comment}{\endcomment}

% For commutative diagrams we use Xy-pic
\usepackage[all]{xy}

% We use 2cell for 2-commutative diagrams.
\xyoption{2cell}
\UseAllTwocells

% We use multicol for the list of chapters between chapters
\usepackage{multicol}

% This is generally recommended for better output
\usepackage{lmodern}
\usepackage[T1]{fontenc}

% For cross-file-references
\usepackage{xr-hyper}

% Package for hypertext links:
\usepackage{hyperref}

% For any local file, say "hello.tex" you want to link to please
% use \externaldocument[hello-]{hello}
\externaldocument[introduction-]{introduction}
\externaldocument[conventions-]{conventions}
\externaldocument[sets-]{sets}
\externaldocument[categories-]{categories}
\externaldocument[topology-]{topology}
\externaldocument[sheaves-]{sheaves}
\externaldocument[sites-]{sites}
\externaldocument[stacks-]{stacks}
\externaldocument[fields-]{fields}
\externaldocument[algebra-]{algebra}
\externaldocument[brauer-]{brauer}
\externaldocument[homology-]{homology}
\externaldocument[derived-]{derived}
\externaldocument[simplicial-]{simplicial}
\externaldocument[more-algebra-]{more-algebra}
\externaldocument[smoothing-]{smoothing}
\externaldocument[modules-]{modules}
\externaldocument[sites-modules-]{sites-modules}
\externaldocument[injectives-]{injectives}
\externaldocument[cohomology-]{cohomology}
\externaldocument[sites-cohomology-]{sites-cohomology}
\externaldocument[dga-]{dga}
\externaldocument[dpa-]{dpa}
\externaldocument[sdga-]{sdga}
\externaldocument[hypercovering-]{hypercovering}
\externaldocument[schemes-]{schemes}
\externaldocument[constructions-]{constructions}
\externaldocument[properties-]{properties}
\externaldocument[morphisms-]{morphisms}
\externaldocument[coherent-]{coherent}
\externaldocument[divisors-]{divisors}
\externaldocument[limits-]{limits}
\externaldocument[varieties-]{varieties}
\externaldocument[topologies-]{topologies}
\externaldocument[descent-]{descent}
\externaldocument[perfect-]{perfect}
\externaldocument[more-morphisms-]{more-morphisms}
\externaldocument[flat-]{flat}
\externaldocument[groupoids-]{groupoids}
\externaldocument[more-groupoids-]{more-groupoids}
\externaldocument[etale-]{etale}
\externaldocument[chow-]{chow}
\externaldocument[intersection-]{intersection}
\externaldocument[pic-]{pic}
\externaldocument[weil-]{weil}
\externaldocument[adequate-]{adequate}
\externaldocument[dualizing-]{dualizing}
\externaldocument[duality-]{duality}
\externaldocument[discriminant-]{discriminant}
\externaldocument[derham-]{derham}
\externaldocument[local-cohomology-]{local-cohomology}
\externaldocument[algebraization-]{algebraization}
\externaldocument[curves-]{curves}
\externaldocument[resolve-]{resolve}
\externaldocument[models-]{models}
\externaldocument[functors-]{functors}
\externaldocument[equiv-]{equiv}
\externaldocument[pione-]{pione}
\externaldocument[etale-cohomology-]{etale-cohomology}
\externaldocument[proetale-]{proetale}
\externaldocument[relative-cycles-]{relative-cycles}
\externaldocument[more-etale-]{more-etale}
\externaldocument[trace-]{trace}
\externaldocument[crystalline-]{crystalline}
\externaldocument[spaces-]{spaces}
\externaldocument[spaces-properties-]{spaces-properties}
\externaldocument[spaces-morphisms-]{spaces-morphisms}
\externaldocument[decent-spaces-]{decent-spaces}
\externaldocument[spaces-cohomology-]{spaces-cohomology}
\externaldocument[spaces-limits-]{spaces-limits}
\externaldocument[spaces-divisors-]{spaces-divisors}
\externaldocument[spaces-over-fields-]{spaces-over-fields}
\externaldocument[spaces-topologies-]{spaces-topologies}
\externaldocument[spaces-descent-]{spaces-descent}
\externaldocument[spaces-perfect-]{spaces-perfect}
\externaldocument[spaces-more-morphisms-]{spaces-more-morphisms}
\externaldocument[spaces-flat-]{spaces-flat}
\externaldocument[spaces-groupoids-]{spaces-groupoids}
\externaldocument[spaces-more-groupoids-]{spaces-more-groupoids}
\externaldocument[bootstrap-]{bootstrap}
\externaldocument[spaces-pushouts-]{spaces-pushouts}
\externaldocument[spaces-chow-]{spaces-chow}
\externaldocument[groupoids-quotients-]{groupoids-quotients}
\externaldocument[spaces-more-cohomology-]{spaces-more-cohomology}
\externaldocument[spaces-simplicial-]{spaces-simplicial}
\externaldocument[spaces-duality-]{spaces-duality}
\externaldocument[formal-spaces-]{formal-spaces}
\externaldocument[restricted-]{restricted}
\externaldocument[spaces-resolve-]{spaces-resolve}
\externaldocument[formal-defos-]{formal-defos}
\externaldocument[defos-]{defos}
\externaldocument[cotangent-]{cotangent}
\externaldocument[examples-defos-]{examples-defos}
\externaldocument[algebraic-]{algebraic}
\externaldocument[examples-stacks-]{examples-stacks}
\externaldocument[stacks-sheaves-]{stacks-sheaves}
\externaldocument[criteria-]{criteria}
\externaldocument[artin-]{artin}
\externaldocument[quot-]{quot}
\externaldocument[stacks-properties-]{stacks-properties}
\externaldocument[stacks-morphisms-]{stacks-morphisms}
\externaldocument[stacks-limits-]{stacks-limits}
\externaldocument[stacks-cohomology-]{stacks-cohomology}
\externaldocument[stacks-perfect-]{stacks-perfect}
\externaldocument[stacks-introduction-]{stacks-introduction}
\externaldocument[stacks-more-morphisms-]{stacks-more-morphisms}
\externaldocument[stacks-geometry-]{stacks-geometry}
\externaldocument[moduli-]{moduli}
\externaldocument[moduli-curves-]{moduli-curves}
\externaldocument[examples-]{examples}
\externaldocument[exercises-]{exercises}
\externaldocument[guide-]{guide}
\externaldocument[desirables-]{desirables}
\externaldocument[coding-]{coding}
\externaldocument[obsolete-]{obsolete}
\externaldocument[fdl-]{fdl}
\externaldocument[index-]{index}

% Theorem environments.
%
\theoremstyle{plain}
\newtheorem{theorem}[subsection]{Theorem}
\newtheorem{proposition}[subsection]{Proposition}
\newtheorem{lemma}[subsection]{Lemma}

\theoremstyle{definition}
\newtheorem{definition}[subsection]{Definition}
\newtheorem{example}[subsection]{Example}
\newtheorem{exercise}[subsection]{Exercise}
\newtheorem{situation}[subsection]{Situation}

\theoremstyle{remark}
\newtheorem{remark}[subsection]{Remark}
\newtheorem{remarks}[subsection]{Remarks}

\numberwithin{equation}{subsection}

% Macros
%
\def\lim{\mathop{\mathrm{lim}}\nolimits}
\def\colim{\mathop{\mathrm{colim}}\nolimits}
\def\Spec{\mathop{\mathrm{Spec}}}
\def\Hom{\mathop{\mathrm{Hom}}\nolimits}
\def\Ext{\mathop{\mathrm{Ext}}\nolimits}
\def\SheafHom{\mathop{\mathcal{H}\!\mathit{om}}\nolimits}
\def\SheafExt{\mathop{\mathcal{E}\!\mathit{xt}}\nolimits}
\def\Sch{\mathit{Sch}}
\def\Mor{\mathop{\mathrm{Mor}}\nolimits}
\def\Ob{\mathop{\mathrm{Ob}}\nolimits}
\def\Sh{\mathop{\mathit{Sh}}\nolimits}
\def\NL{\mathop{N\!L}\nolimits}
\def\CH{\mathop{\mathrm{CH}}\nolimits}
\def\proetale{{pro\text{-}\acute{e}tale}}
\def\etale{{\acute{e}tale}}
\def\QCoh{\mathit{QCoh}}
\def\Ker{\mathop{\mathrm{Ker}}}
\def\Im{\mathop{\mathrm{Im}}}
\def\Coker{\mathop{\mathrm{Coker}}}
\def\Coim{\mathop{\mathrm{Coim}}}

% Boxtimes
%
\DeclareMathSymbol{\boxtimes}{\mathbin}{AMSa}{"02}

%
% Macros for moduli stacks/spaces
%
\def\QCohstack{\mathcal{QC}\!\mathit{oh}}
\def\Cohstack{\mathcal{C}\!\mathit{oh}}
\def\Spacesstack{\mathcal{S}\!\mathit{paces}}
\def\Quotfunctor{\mathrm{Quot}}
\def\Hilbfunctor{\mathrm{Hilb}}
\def\Curvesstack{\mathcal{C}\!\mathit{urves}}
\def\Polarizedstack{\mathcal{P}\!\mathit{olarized}}
\def\Complexesstack{\mathcal{C}\!\mathit{omplexes}}
% \Pic is the operator that assigns to X its picard group, usage \Pic(X)
% \Picardstack_{X/B} denotes the Picard stack of X over B
% \Picardfunctor_{X/B} denotes the Picard functor of X over B
\def\Pic{\mathop{\mathrm{Pic}}\nolimits}
\def\Picardstack{\mathcal{P}\!\mathit{ic}}
\def\Picardfunctor{\mathrm{Pic}}
\def\Deformationcategory{\mathcal{D}\!\mathit{ef}}
\usepackage{cleveref}
\crefname{originaldefinition}{original definition}{original definitions}
\Crefname{originaldefinition}{Original definition}{Original definitions}
\makeatletter
\newlabel{situation-complexes}{{0DLE}{source}{}{sourcepointer-1}{}}
\newlabel{situation-complexes@cref}{{[originaldefinition][0][]0DLE}{[1][0][]source}{}{}{}}
\makeatother
\crefname{originaldefinition}{original definition}{original definitions}
\Crefname{originaldefinition}{Original definition}{Original definitions}
\makeatletter
\newlabel{artin-definition-limit-preserving}{{07XL}{source}{}{sourcepointer-2}{}}
\newlabel{artin-definition-limit-preserving@cref}{{[originaldefinition][0][]07XL}{[1][0][]source}{}{}{}}
\makeatother
\crefname{originalremark}{original remark}{original remarks}
\Crefname{originalremark}{Original remark}{Original remarks}
\makeatletter
\newlabel{spaces-limits-remark-limit-preserving}{{05N0}{source}{}{sourcepointer-3}{}}
\newlabel{spaces-limits-remark-limit-preserving@cref}{{[originalremark][0][]05N0}{[1][0][]source}{}{}{}}
\makeatother
\crefname{originallemma}{original lemma}{original lemmas}
\Crefname{originallemma}{Original lemma}{Original lemmas}
\makeatletter
\newlabel{spaces-more-morphisms-lemma-descend-relatively-perfect}{{0DKW}{source}{}{sourcepointer-4}{}}
\newlabel{spaces-more-morphisms-lemma-descend-relatively-perfect@cref}{{[originallemma][0][]0DKW}{[1][0][]source}{}{}{}}
\makeatother
\crefname{originallemma}{original lemma}{original lemmas}
\Crefname{originallemma}{Original lemma}{Original lemmas}
\makeatletter
\newlabel{lemma-complexes}{{0DLD}{source}{}{sourcepointer-5}{}}
\newlabel{lemma-complexes@cref}{{[originallemma][0][]0DLD}{[1][0][]source}{}{}{}}
\makeatother
\crefname{originallemma}{original lemma}{original lemmas}
\Crefname{originallemma}{Original lemma}{Original lemmas}
\makeatletter
\newlabel{lemma-complexes-open-neg-exts-vanishing}{{0DLC}{source}{}{sourcepointer-6}{}}
\newlabel{lemma-complexes-open-neg-exts-vanishing@cref}{{[originallemma][0][]0DLC}{[1][0][]source}{}{}{}}
\makeatother
\crefname{originallemma}{original lemma}{original lemmas}
\Crefname{originallemma}{Original lemma}{Original lemmas}
\makeatletter
\newlabel{limits-lemma-limit-contained-in-constructible}{{05F4}{source}{}{sourcepointer-7}{}}
\newlabel{limits-lemma-limit-contained-in-constructible@cref}{{[originallemma][0][]05F4}{[1][0][]source}{}{}{}}
\makeatother
\title{The stack of relatively perfect complexes is limit preserving}
\author{The Stacks Project Authors}
\date{stacks-project:quot.tex:local-snapshot-20261012; GNU FDL 1.2 or later}
\begin{document}
\maketitle
\setcounter{subsection}{1}
\noindent\textit{Editorial note.} All delivered mathematical text is the author's own, copied from the pinned Stacks Project chapter file quot.tex; no statement or proof step was rewritten, reordered or removed. The counted span is the author environment and its own proof environment. References inside the span are handed to the assembler as navigation stubs that carry the canonical tag of the cited tag only; whether each of them is load-bearing is recorded on the stub itself.\par\medskip
\noindent\textit{Allowed standard prerequisites (editorial).}\begin{itemize}
\item Situation situation-complexes (canonical tag 0DLE). cited at statement line 5614. Standing situation: X -> B and the relatively perfect K.
\item Definition artin-definition-limit-preserving (canonical tag 07XL). cited at statement line 5617. The property being established: limit preservation.
\item Remark spaces-limits-remark-limit-preserving (canonical tag 05N0). cited at proof line 5636. Bottom arrow equivalence; reduce to limits over B.
\item Lemma spaces-more-morphisms-lemma-descend-relatively-perfect (canonical tag 0DKW). cited at proof lines 5642-5643. T-perfect objects = colimit of T\_i-perfect objects.
\item Lemma lemma-complexes (canonical tag 0DLD). cited at proof lines 5650 and 5653. Condition (2) being descended.
\item Lemma lemma-complexes-open-neg-exts-vanishing (canonical tag 0DLC). cited at proof line 5655. The open W of negative-Ext vanishing for E\_i.
\item Lemma limits-lemma-limit-contained-in-constructible (canonical tag 05F4). cited at proof line 5660. T -> T\_i factors through W at a finite stage i'.
\end{itemize}\medskip
\ifdefined\hypertarget\hypertarget{sourcepointer-1}{}\fi
\noindent\textit{Source reference.} Original definition 0DLE: sources/stacks-chow/tags.txt, source lines 17044--17044. Cited at statement line 5614: the lemma is phrased in this situation; load-bearing.\par\medskip
\ifdefined\hypertarget\hypertarget{sourcepointer-2}{}\fi
\noindent\textit{Source reference.} Original definition 07XL: sources/stacks-chow/tags.txt, source lines 10081--10081. Cited at statement line 5617: the property being established; load-bearing.\par\medskip
\ifdefined\hypertarget\hypertarget{sourcepointer-3}{}\fi
\noindent\textit{Source reference.} Original remark 05N0: sources/stacks-chow/tags.txt, source lines 7289--7289. Cited at proof line 5636: the bottom arrow of the diagram is an equivalence since B -> S is locally of finite presentation; load-bearing.\par\medskip
\ifdefined\hypertarget\hypertarget{sourcepointer-4}{}\fi
\noindent\textit{Source reference.} Original lemma 0DKW: sources/stacks-chow/tags.txt, source lines 17026--17026. Cited at proof lines 5642-5643: the category of T-perfect objects is the colimit; load-bearing.\par\medskip
\ifdefined\hypertarget\hypertarget{sourcepointer-5}{}\fi
\noindent\textit{Source reference.} Original lemma 0DLD: sources/stacks-chow/tags.txt, source lines 17043--17043. Cited at proof lines 5650 and 5653: condition (2) of that lemma; load-bearing twice.\par\medskip
\ifdefined\hypertarget\hypertarget{sourcepointer-6}{}\fi
\noindent\textit{Source reference.} Original lemma 0DLC: sources/stacks-chow/tags.txt, source lines 17042--17042. Cited at proof line 5655: the open W constructed for E\_i; load-bearing.\par\medskip
\ifdefined\hypertarget\hypertarget{sourcepointer-7}{}\fi
\noindent\textit{Source reference.} Original lemma 05F4: sources/stacks-chow/tags.txt, source lines 7013--7013. Cited at proof line 5660: T -> T\_i factors through W at some finite stage i'; load-bearing.\par\medskip
\begin{lemma}
\label{lemma-complexes-limits}
In Situation \ref{situation-complexes} assume that $B \to S$
is locally of finite presentation. Then
$p : \Complexesstack_{X/B} \to (\Sch/S)_{fppf}$ is limit preserving
(Artin's Axioms, Definition \ref{artin-definition-limit-preserving}).
\end{lemma}

\begin{proof}
Write $B(T)$ for the discrete category whose
objects are the $S$-morphisms $T \to B$. Let $T = \lim T_i$ be a filtered
limit of affine schemes over $S$. Assigning to an object
$(T, h, E)$ of $\Complexesstack_{X/B, T}$ the object $h$
of $B(T)$ gives us a commutative diagram of fibre categories
$$
\xymatrix{
\colim \Complexesstack_{X/B, T_i} \ar[r] \ar[d] &
\Complexesstack_{X/B, T} \ar[d] \\
\colim B(T_i) \ar[r] & B(T)
}
$$
We have to show the top horizontal arrow is an equivalence. Since
we have assume that $B$ is locally of finite presentation over $S$
we see from
Limits of Spaces, Remark \ref{spaces-limits-remark-limit-preserving}
that the bottom horizontal arrow is an equivalence. This means that
we may assume $T = \lim T_i$ be a filtered limit of affine schemes over
$B$. Denote $g_i : T_i \to B$ and $g : T \to B$ the corresponding
morphisms. Set $X_i = T_i \times_{g_i, B} X$ and $X_T = T \times_{g, B} X$.
Observe that $X_T = \colim X_i$.
By More on Morphisms of Spaces, Lemma
\ref{spaces-more-morphisms-lemma-descend-relatively-perfect}
the category of $T$-perfect objects of $D(\mathcal{O}_{X_T})$
is the colimit of the categories of $T_i$-perfect objects
of $D(\mathcal{O}_{X_{T_i}})$.
Thus all we have to prove is that given an $T_i$-perfect object
$E_i$ of $D(\mathcal{O}_{X_{T_i}})$ such that
the derived pullback $E$ of $E_i$ to $X_T$ satisfies
condition (2) of Lemma \ref{lemma-complexes},
then after increasing $i$ we have that
$E_i$ satisfies
condition (2) of Lemma \ref{lemma-complexes}.
Let $W \subset |T_i|$ be the open constructed
in Lemma \ref{lemma-complexes-open-neg-exts-vanishing}
for $E_i$ and $E_i$. By assumption on $E$ we find
that $T \to T_i$ factors through $T$.
Hence there is an $i' \geq i$
such that $T_{i'} \to T_i$ factors through $W$, see
Limits, Lemma \ref{limits-lemma-limit-contained-in-constructible}
Then $i'$ works by construction of $W$.
\end{proof}

\end{document}
